\scanpage{8}\mspage{7}

On peut écrire

\[
D=\sum \unsure{e_i|e^i}\otimes e'_i
\]

où $(e_i)$ est une base de $H^1(X)$, et $e'_i$ est la base duale de $H^1(B)$. Posant

\[
\mathrm{Todd}_{X/k}\exp\Delta=\sum_{i=0}^{n}a^i\qquad a^i\in H^{2i}(X)
\]

on trouve

\[
\exp D\,\mathrm{pr}_1^*(\mathrm{Todd}_{X/k}\exp\Delta)=\sum_{i,j}\mathrm{pr}_1^*(a^i)\,D^j/j!
\]

le terme en $i,j$ écrit étant de degré cohomologique partiel en le facteur $X$ égal à $2i+j$. Il donne une contribution nulle dans $\mathrm{pr}_{2*}$ sauf si $2i+j=2n$ i.e. $j=2(n-i)$, \deleted{donc il faut} \inserted{auquel cas on trouve} un $\mathrm{pr}_{2*}$ de degré $j$, qui n'est intéressant que si $j=2$ i.e. $n-i=1$. \deleted{prendre} Donc on trouve

\[
\deleted{\unsure{\varphi(\delta)=|\sum}\tfrac12\mathrm{pr}_1^*(a_{n-1})D^2\unsure{\in H^2(B)}.}
\]

\[
\varphi(\delta)\deleted{\unsure{\in}}=\tfrac12\mathrm{pr}_{2*}\bigl(\mathrm{pr}_1^*(a_{n-1})D^2\bigr)
\]

\uline{Ex. 1}

Prenons le cas $n=1$. Alors $\delta$ est caractérisé par son degré, soit \deleted{\unsure{$\delta$}} $d$. On trouve

\[
\deleted{\mathrm{Todd}_{X/k}\exp\Delta=(1\unsure{-|+}\tfrac K2)(1+\Delta)=1+\underbrace{(\Delta-\tfrac K2)}_{\substack{\text{degré}\\d-g+1}}\unsure{\nearrow|\diagup}}
\]

\[
\varphi(\delta)=\tfrac12\mathrm{pr}_{2*}(D^2)
\]

\deleted{et il faut prendre $\mathrm{pr}_{2*}$ de}

\[
\deleted{\tfrac12 D^2+\mathrm{pr}_1^*(\Delta-\tfrac K2)}
\]

Or \deleted{\unsure{nous}} on a, si $(e_1,f_1,\unsure{e_2,f_2|e,f},\ldots,e_g,f_g)$ est une base symplectique de $H^1(X)$ [$e_i f_i=-f_i e_i=\eta$, base can. de $H^2(X)$),

% slanted
\margin{On identifie $H^1(X)$ à son dual grâce à la forme \unsure{intersection|d'intersection}.}

\[
D=\sum(e_i\otimes f_i+f_i\otimes e_i)
\]
